\chapter{保护分区分析与模板验证}

\section{保护分区流程}

保护分区分析通常需要把生态重要性和管理可行性放在同一框架内。示例流程先根据物种丰富度和栖息地质量识别生态热点，再结合连通性指标判断廊道价值，最后叠加人为扰动和管理边界形成分区建议。该流程并不代表真实规划结论，仅用于展示博士论文中常见的方法叙述。

\begin{enumerate}
  \item 整理物种记录和环境变量，生成统一分析网格；
  \item 计算物种丰富度、栖息地质量和廊道连通性；
  \item 叠加人为扰动、道路距离和保护地边界；
  \item 按生态重要性和管理紧迫性划分优先等级；
  \item 对分区结果进行不确定性解释和敏感性分析。
\end{enumerate}

\section{示例结果}

为了展示多列表格排版，表 \ref{tab:nwu-zone-result} 给出一组公开演示数据。数据仅用于模板验证，不对应真实监测结果。

\begin{table}[htbp]
  \centering
  \caption{保护分区指标示例}
  \label{tab:nwu-zone-result}
  \begin{tabular}{cccccc}
    \toprule
    区域 & 丰富度指数 & 连通性指数 & 扰动压力 & 不确定性 & 优先等级\\
    \midrule
    A & 0.82 & 0.76 & 0.21 & 0.18 & 高\\
    B & 0.68 & 0.63 & 0.32 & 0.22 & 中高\\
    C & 0.54 & 0.48 & 0.44 & 0.25 & 中\\
    D & 0.37 & 0.41 & 0.58 & 0.31 & 一般\\
    \bottomrule
  \end{tabular}
\end{table}

示例结果表明，模板能够稳定呈现章节标题、正文段落、枚举列表、表格、图题和公式。正式论文还应补充数据来源、样本量、模型参数、误差分析、敏感性检验和管理建议。

\section{模板实现说明}

`thesis-nwu-doctor` 在 `packages/bensz-thesis/` 中注册独立 profile 和 style，项目层只维护元信息、公开正文和项目说明。这样的分层能够避免修改西北大学博士论文模板时影响南京大学、中国科学院大学、中山大学、哈尔滨工业大学等既有模板。
